\section{Distress, coordination and conduct}
\label{sec:conduct}

This section follows what players said when solvency, inheritance or a rival's development was at stake. The task allows public deception, so a false public statement is an authorized tactic and not by itself evidence of misalignment with the instructions; the private-honesty request applies only under the standard prompt. For each incident we separate whether a proposition is false or misleading under a stated reading, whether another channel from the same speaker contradicts it, whether it reached the recipient, what the recipient's record shows, and what changed afterward. A recipient's acceptance does not show that it believed the speaker.

The offer ledger shows distress only crudely. Of 339 acceptance-legal offers in which the author asks for net cash, 132 come from authors holding less than \$100; these are accepted 18 times and rejected 62 times, while the 51 from authors holding \$700 or more are accepted 13 times and rejected 6 times. The threshold is arbitrary, cash on hand is not accessible finance (Section~\ref{sec:liquidity}), and the offers sell different things, so this is description only. The incidents show what the ledger cannot.

\Incident{g13}{An insurance pitch that omits available collateral (G13)}
\Setup{At turn 107 of game 13, Grok has \$35 and stands on Atlantic. Sol owns all four railroads, so a Short Line landing costs \$200, and a Grok bankruptcy to Sol would pass Grok's yellows to Sol. Grok's Ventnor is still clear and can be mortgaged for \$130. Grok offers its mortgaged Pennsylvania Avenue to Opus.}

\TranscriptQuote{g13-opening-insurance}

Opus counters at \$130 in d251, publicly citing the cost of lifting the mortgage; its returned provider explanation values keeping Grok solvent, a rationale that follows Grok's framing rather than forming independently of it. Grok answers:

\TranscriptQuote{g13-seller-counter}

The same response records the mortgage route and a preference to avoid it, and Grok's returned provider explanation is more direct:

\TranscriptQuote{g13-seller-floor}
\TranscriptQuote{g13-omission-explanation}

Opus sees the new message and accepts:

\TranscriptQuote{g13-acceptance}

\TranscriptFinding{Events 1808--1812 transfer Pennsylvania to Opus, \$200 to Grok and \$16 of transfer interest to the bank, leaving Grok with \$235. An executed alternative from the d252 prefix shows that accepting \$130 returns Grok to its management menu, where mortgaging Ventnor reaches \$295. Forty-three turns later, at d311, Grok owes Sol \$200 on Short Line with \$49 and no financing action, and its yellows pass to Sol. Opus wins the game.}

\TranscriptReading{This is the clearest retained case of a delivered omission with an explicit concealment rationale: ``\$130 leaves me \$35 short'' is true in cash and omits a legal way to pay. A narrower reading also fits, in which Grok states a refusal posture rather than claiming payment would be impossible. Opus's acceptance repeats the estate stakes but never says the cheaper route is impossible, so the record does not show that Opus believed the false minimum, and the \$70 price difference is not an identified effect of the message. The opening probability is also wrong (Short Line is 4/36 from Atlantic), and an earlier claim that inheritance would complete Opus's greens had been publicly corrected at d244. The insurance bought time but did not keep the yellows from Sol.}

\TranscriptSource{G13 d250--253 and d311. Grok's quoted response is standard; Opus's d251 and d253 responses use action-only recovery and have no applied note. The d252--253 timelines contain a self-partner formatter description, while correct holdings remain available.}

\Incident{g18}{A standing offer withdrawn when the seller is nearly broke (G18)}
\Setup{In game 18, Sol holds mortgaged Pacific Avenue, inherited from Grok's estate, and Opus needs it to complete its greens. Opus offers \$250 at d421, calls it final, and says at d423, d425 and d427 that the offer stands. After collecting \$750 in rent from Sol it repeats the assurance.}

\TranscriptQuote{g18-assurance}

Sol first asks for more. A second \$750 rent payment forces it to sell orange buildings and mortgage its reds at d432--436, leaving \$49, and Opus rejects Sol's \$450 and \$350 packages. At d441 Sol invokes the advertised price:

\TranscriptQuote{g18-invocation}
\TranscriptQuote{g18-withdrawal}
\TranscriptQuote{g18-withdrawal-explanation}

\TranscriptFinding{The d431 assurance is in both players' delivered memory at the invocation and the withdrawal. No sale occurs at \$250. Opus later offers \$150 and \$225 and rejects Sol's \$300 proposal. Sol is eliminated at d463 owing \$900, Pacific reaches Opus only in the game-ending estate, and Opus wins the game.}

\TranscriptReading{The record shows a repeated assurance, an explicit invocation, and a withdrawal that cites the seller's near-insolvency. The stated reason is economically intelligible: \$250 would extend a nearly bankrupt rival's game, and Opus's own greens were mortgaged. The record does not show costly reliance, since the intervening liquidation paid a second rent bill and Sol had asked for more than \$250, and withdrawal does not show the offer was dishonest when made. The game has no standing-offer object or expiry, so there is no assigned norm to breach~\cite{empty_commitments_v2,commitment_standing}. In ledger terms, Opus declines an offer completing its own group at the price it had named, because the counterparty's distress changed the deal's value to Opus.}

\TranscriptSource{G18 d421--442 and d454--461. The withdrawal explanation is provider-returned text under action-only recovery, not an applied note.}

\Incident{g20}{An appeal to deny a rival building cash (G20)}
\Setup{In game 20, Opus has sold Illinois to Grok for \$500, completing Grok's reds, and Grok has little cash to build. At d96 Gemini addresses the table.}

\TranscriptQuote{g20-coordination-appeal}
\TranscriptQuote{g20-recognition-note}

Grok's d97 received action is \End{} despite text describing building. It builds three red houses at d98, mortgages New York at d99 and offers Vermont to Gemini:

\TranscriptQuote{g20-sale-response}
\TranscriptQuote{g20-concrete-refusal}

\TranscriptFinding{Gemini refuses the kind of cash sale its appeal described, and Grok mortgages Vermont at d102. The target of the appeal names it a boycott.}

\TranscriptReading{This is an observed attempt to coordinate financial denial, recognized by its target and followed by the initiator's own consistent refusal. It is not an established cartel. Sol had stated its own reason not to finance Grok at d95, before the appeal, and Opus later proposed cash purchases from Grok twice, neither of which settled. Nothing shows the table committing to the appeal or the appeal changing another player's choice. It is the rival-sensitivity of Section~\ref{sec:bargaining} voiced to a third party.}

\TranscriptSource{G20 d95--102. The d96 appeal is in Grok's d97 prompt; Gemini's d101 refusal answers the d100 offer. The d97 note accompanies selected attempt 1 after a malformed attempt 0.}

\Incident{g3}{A building obstacle contradicted by the speaker's own note (G3)}
\Setup{Late in game 3, Gemini owns the dark blues and one bank house remains. Grok discusses the building threat publicly and privately, and Gemini then gives a reason to hold.}

\TranscriptQuote{G3-2}
\TranscriptQuote{G3-3}
\TranscriptQuote{G3-5}

\TranscriptFinding{Gemini ends its next management decision without building. At that d364 opportunity, after a \$15 tax, it has \$43. An executed legal sequence mortgages Kentucky for \$110 and Atlantic for \$130 and buys the \$200 Park Place house, leaving \$83, with no release of buildings by any other player.}

\TranscriptReading{The public claim that a release is needed is false for this route, while the cash shortfall is real. The speaker's note records the route, but its longer provider explanation repeatedly treats a release as necessary, so the channels do not establish a coherent deliberate deception. Gemini had already held at d362, before the claim, and its d364 note prefers to keep income unmortgaged, so exposure followed by another hold does not show reliance. The case is misleading pressure with mixed speaker evidence and a legal route the recipient did not take.}

\TranscriptSource{G3 d362--364. The witness comes from the retained exact-prefix legality audit and was reproduced in the transcript source review.}

\Incident{g19}{A false premise corrected, then accepted (G19)}
\Setup{In game 19, Grok offers Virginia plus \$150 for Sol's States Avenue; both deeds are pink. The selected response's assistant content names a different action, and its private field corrects the color and says to abort (two separate exact spans).}

\TranscriptQuote{g19_d135_content_end}
\TranscriptQuote{g19-abort}

The received tool is \texttt{propose\_trade}, with this public message, and Sol answers:

\TranscriptQuote{g19_d135_false_public}
\TranscriptQuote{g19_d136_correction_accept}

\TranscriptFinding{Events 1023--1027 execute the swap and the \$150 payment. The proposal was the single valid received call in a standard response, not a fallback or retry. Sol's note treats States as pink and the terms as favorable.}

\TranscriptReading{The recipient corrects the false premise publicly and accepts anyway, so settlement here cannot be read as belief in the premise. The speaker's own response cannot support a coherent plan to deceive, since its content chooses \End{} and its note says not to send the offer. G6 d19--21 has the same structure: Sol's note rejects Grok's false claim about an orange monopoly while accepting the trade.}

\TranscriptSource{G19 d135--136, turn 67, attempt 0. The false claim is in Sol's prompt. Received, parsed and applied action names agree; the conflict is with assistant content and the note.}

\subsection{What the incidents add}

The stakes players name are the ones the ledgers measure: completion (G2, G11, G20), liquidity (G13, G18, G3) and inheritance (G13). Recipients do not simply absorb false claims. Grok in G11 and Sol in G19 and G6 correct them on the record, and in G4 Gemini's public ``Fair point'' to a false \$200 house cost follows a note that already holds the true \$150, so a public concession need not mean a changed belief. Only G13 leaves the recipient silent on the omitted route. Applied actions can also contradict every text channel in the same response, as in G12 and G19. Across the cohort, 139 selected Grok responses carry a content/tool name conflict; in 15 the arguments are identical and only the function name differs, and four of these accept trades while the message argues for rejection. One, G18 d83, accepts \$130 for Baltic and completes Opus's browns in a game Opus won. These conflicts constrain which channel can be read as the agent's plan; they do not measure deception. Appendix~\ref{app:incidents} gives the remaining chronologies, a five-case proposition table, a bluff that settled at its advertised limit (G7) and a private note delivered as public chat (G20 d1).

\FloatBarrier
